\section{Introduction}

Revealing the paths of selected cards can separate a shuffle's remaining
randomness into local choices, but it also changes the law of those choices.
In an ordered coordinate sweep, conditioning still separates by coordinate
lines; each residual choice is a bijection between the slots left free by the
prescribed paths at that stage. This paper controls the entire remaining bijection while
retaining the exact cost of prescribing distinct positions without
replacement. That cost makes the estimate stable when further card paths
are revealed during the proof.

Fix a positive integer \(d\) and put \(N=2^d\). A binary coordinate sweep
acts on the slots \(\{0,1\}^d\). In each coordinate direction it
independently fixes or swaps the two cards on every parallel edge, with
probability \(1/2\) each, and visits the coordinates in order. Write \(B_N\)
for its law. One card is uniform after a sweep, but the joint permutation
retains dependence because cards can use the same switches.

The physical model comes from Thorp's study of nonrandom shuffling and
Faro~\cite{Thorp1973}. Split a deck into equal halves, pair corresponding
cards, and use an independent fair coin to order each pair before
interleaving. On binary position strings one physical Thorp shuffle is
\[
 (x_1,x_2,\ldots,x_d)\longmapsto
 (x_2,\ldots,x_d,x_1+\xi_{x_2,\ldots,x_d}),
\]
where addition is modulo two and the bits \(\xi\) are independent and
fair. Undoing the deterministic coordinate rotation makes successive
switch layers act in successive coordinate directions. After \(d\)
physical shuffles the rotation is the identity, so one binary sweep costs
exactly \(d\) physical steps; see also
Morris~\cite[Section~1.1]{Morris2008} for the coordinate representation.

Let \(q_d\) be the law of one physical step and \(U_N\) the uniform
probability on \(S_N\). We use
\[
 \|\mu-\nu\|_{\TV}
 =\frac12\sum_{g\in S_N}|\mu(g)-\nu(g)|,\qquad
 t_{\rm mix}(d)=\inf\{t\in\mathbb Z_{\ge0}:
                 \|q_d^{*t}-U_N\|_{\TV}\le1/4\}.
\]
Translation by a deterministic initial deck preserves the distance to
uniform, so this definition is the same for every deterministic initial
deck.

For the irreducible unitary representation \(\rho_\lambda\) of \(S_N\)
indexed by \(\lambda\vdash N\), put
\[
 K_\lambda=\mathbb E_{g\sim B_N}\rho_\lambda(g),
 \qquad D_\lambda=\dim\rho_\lambda.
\]
The binary consequence of our argument is the following dimension bound.

\begin{theorem}[Binary sweep contraction]\label{thm:binary}
There are absolute constants \(g>0\) and \(d_0\) such that, for \(d\ge d_0\)
and every \(\lambda\vdash 2^d\),
\[
 \|K_\lambda\|_{\op}\le D_\lambda^{-g}.
\]
The sign representation is annihilated. There is an absolute positive integer \(w\)
such that the law after \(w\) independent sweeps converges to uniform in
total variation as \(d\to\infty\), uniformly over deterministic initial
decks.
\end{theorem}

The number of random bits supplies a matching lower bound in order.

\begin{proposition}[Support obstruction]\label{prop:support}
For every integer \(t\ge0\),
\[
 \|q_d^{*t}-U_N\|_{\TV}\ge1-\frac{2^{tN/2}}{N!}.
\]
Consequently,
\[
 t_{\rm mix}(d)\ge
 \left\lceil\frac2N\log_2\!\left(\frac{3N!}{4}\right)\right\rceil
 =2d-O(1).
\]
\end{proposition}
\begin{proof}
The \(tN/2\) fair bits determine the permutation, whose support therefore
has at most \(2^{tN/2}\) elements. Testing that support in the definition
of total variation gives the first bound. Distance at most \(1/4\)
requires its uniform mass to be at least \(3/4\), proving the integer
bound. Stirling's formula gives the asymptotic expression.
\end{proof}

Theorem~\ref{thm:binary} and Proposition~\ref{prop:support} give
\(t_{\rm mix}(d)=\Theta(d)\) in physical shuffle steps. In each fixed
dimension, convergence also holds. A sweep can be the identity or any
single cube-edge transposition with positive probability, and the edge
transpositions of a connected graph generate \(S_N\). Padding products
with identities therefore gives a power of the sweep law with full
support. Its positive minimum probability gives a uniform minorization
and convergence by iteration.

\subsection{Earlier work}

Morris's July 2005 author version, published in journal form in 2008,
gave the polynomial upper bound \(O(d^{44})\) physical shuffles for
\(N=2^d\) cards~\cite{Morris2008}. Montenegro and Tetali sharpened this
to \(O(d^{29})\)~\cite{MontenegroTetali2006}. Morris then obtained
\(O((\log N)^4)\) physical shuffles for every even deck size
\(N\)~\cite{Morris2009}, followed by \(O(d^3)\) for power-of-two
decks~\cite{Morris2013}.

Earlier analyses already use information revealed about card motion.
Morris's chameleon identity represents a tagged-card law conditioned on
the trajectory of an unordered occupied set
\cite[Section~4, Lemma~3]{Morris2008}. In his later entropy proof,
complete labeled trajectories in two coupled runs of the reverse Thorp shuffle are
exposed successively, and the expected entropy changes are summed to
obtain contraction after \(d\) physical shuffles
\cite[Section~4, proof of Lemma~7]{Morris2013}.
Morris, Rogaway and Stegers analyze selected-card Thorp marginals by
averaged conditional one-card bounds and a coupling along the evolution
\cite[Section~3 and Appendix~A, Lemma~2]{mrs}. Hoang, Morris and Rogaway
use a related conditional squared-discrepancy calculation for the
distinct swap-or-not shuffle
\cite[Section~3, proof of Theorem~3 and Lemma~4]{hmr2014}.

Czumaj and V\"ocking also studied partial Thorp permutations through
non-Markovian coupling. For dyadic decks and each fixed
\(0<\varepsilon<1\), their result gives nearly uniform joint positions
for a fraction \(1-\varepsilon\) of the labeled cards after \(O(d^2)\) physical shuffles
\cite{CzumajVocking2014}; equivalently, \(O(d)\) complete coordinate sweeps
\cite[Section~1.5]{Czumaj2015}. That result concerns the endpoint law of
selected cards after repeated binary sweeps. Our conditional theorem
concerns the entire residual bijection for each feasible fixed family
of paths under line laws close to uniform.

The representation ingredients are classical branching, Pieri's rule,
the hook-length formula, and ordinary Schur--Weyl duality
\cite{Sagan2001,Stanley1999}. The signed tensor action
and its hook support belong to the hook theory of Berele and
Regev~\cite{BereleRegev1983}. Positive domination by mixtures of tensor
powers also underlies the postselection method of Christandl, K\"onig and
Renner~\cite{ChristandlKonigRenner2009}. We derive the quantitative signed
projection estimate needed here from the ordinary Schur--Weyl sectors.
The analytic ingredient is the classical subharmonic maximum principle;
Section~\ref{sec:analytic} gives the local barrier argument.

\subsection{The conditional estimate and the proof}

The main conditional input is Theorem~\ref{thm:conditional}. It treats an
ordered product grid whose coordinate sizes lie in a fixed interval
\([R,R^2]\) of powers of two. On a coordinate line of size \(m\), the law
is a mixture \((1-z)U_m+zB_m\) of a uniform permutation and a binary
sweep. The theorem holds on a small real interval \(0\le z\le z_*\),
where these line laws are uniformly comparable to \(U_m\).

In one ordered sweep, a card's input and output determine its whole path:
after stage \(j\), the first \(j\) coordinates have their output values
and the remaining coordinates still have their input values. On the small
real interval, a prescribed family is feasible exactly when these forced paths occupy distinct slots
at every layer. The conditional theorem bounds a Schatten moment of the
remaining random bijection, in every irreducible representation, for
each such family. Its bound has a negative term proportional to the
logarithm of the representation dimension, an allowance for the number
of prescribed cards, and a negative
correction for sampling their positions without replacement. On a line
of size \(m\) with \(u\) prescribed assignments, this correction is
\(\log(m^u/(m)_u)\), where \((m)_u=m(m-1)\cdots(m-u+1)\) and \((m)_0=1\).
For a prescribed family \(\mathcal H\), let \(C(\mathcal H)\) be the sum
of these corrections over all stage lines.

The correction is the invariant that permits more paths to be revealed.
When a placement of additional cards augments \(\mathcal H\) to
\(\mathcal H^+\), its transition probability is bounded using the difference
\(C(\mathcal H^+)-C(\mathcal H)\) of the summed line corrections.
Later, branching realizes the removal of boxes from a Young diagram by
recording such additional placements. A cyclic trace expansion and
H\"older's inequality use that difference to return the original
\(-C(\mathcal H)\) after the placements are summed. Because these fibers
carry augmented path constraints, the induction must hold simultaneously
for the conditioned laws.

The proof first treats sparse representation levels: types that first appear
when only a short list of cards is tracked, together with their sign
twists. Alternating the placement probabilities over subsets of the list
cancels any tracked path whose stage lines are untouched by all other
tracked or prescribed paths. With many coordinates, a count using rooted
forests makes it unlikely that every tracked path shares a line. With a
bounded number of coordinates, lines are larger and this count is too
weak. At \(z=0\), an inverse-gamma moment identity expresses the
falling-factorial ratios as expectations. A tracked path may now share
lines: if it shares no layer slot with another tracked or prescribed
path, and the number of tracked and prescribed paths using each of its
lines is a small fraction of that line's size, alternation produces small
differences of the line factors.
Factors from particles sharing a line are grouped before their moments
are bounded. Once \(R\) is fixed, only finitely many grids have this
bounded number of coordinates, so continuity extends the estimate to a
common small real interval.

To handle the remaining types of large dimension, choose \(p\) and first
prove an auxiliary estimate for every diagram whose boxes lie in the
first \(p\) rows or the first \(p\) columns, uniformly over all admitted
path families on the current grid. Split the coordinate list into two
blocks and regard the grid
as a rectangle: the first block sweeps within rows, and the second
within columns. At their junction the prescribed paths have removed
some sites. In the signed tensor representation, each row or column
isotypic projection is bounded by a controlled multiple of a probability
mixture of tensor powers of positive trace-one matrices. Averaging the
chosen row matrices supplies a trace-one reference matrix for bounding
the overlap of the row and column projections on the free sites.
The child moment estimates then cancel the cost of the subgroup dimensions in
that overlap, leaving a strengthened estimate for the full hook type
with the same trajectory correction. Finally, adding the boxes outside
the hook creates the additional path fibers described above. Branching
supplies a factor equal to the dimension of the removed diagram; together
with the remaining dimension saving, this pays for the extra paths.

We then set \(\mathcal H=\varnothing\). The unconditioned sweep operator
is a matrix polynomial in \(z\), and the conditional theorem gives a
dimension saving on a short real segment. For the finitely many allowed
line sizes, the binary line averages have a strict norm gap away from
their invariant vectors. This bounds the full polynomial operator by
one on a complex disk of radius greater than one. Applying the
subharmonic maximum principle to scalar matrix coefficients on the disk
with the segment removed transfers the dimension saving to \(z=1\),
where the sweep is binary. Finite-group Fourier analysis then gives the
fixed number of sweeps in Theorem~\ref{thm:binary}.

Section~\ref{sec:prelim} fixes the representation conventions.
Section~\ref{sec:paths} defines the conditional law and states its moment
bound. Section~\ref{sec:sparse} proves the sparse estimate, and
Section~\ref{sec:density} constructs positive tensor densities for the
signed action. Section~\ref{sec:induction} combines these inputs and
completes the moment induction by adding trajectory constraints.
Section~\ref{sec:analytic} carries the
unconditioned estimate to the binary sweep and proves the mixing
consequence.
